\newpage

\KOMAoptions{paper=A3,paper=landscape}
\areaset{390mm}{270mm}
\fancyheadoffset{0pt} % обновляем хедер и футер

\topmargin=-2.6cm % Команда \topmargin задаёт верхнее поле страницы. При этом поле отсчитывается не от левого края листа, а от линии, параллельной краю листа и отстоящей от него на 1 дюйм
\oddsidemargin=0cm % отступ слева https://zns.susu.ru/IT/latex/Marking/marking.html
\headsep=0cm
\footskip=-0.8cm
\setlength{\headheight}{1cm}

\color{uniblue}{\section[(рекомендуемое) Структурно-функциональная схема устройства]{(рекомендуемое)\\Структурно-функциональная схема устройства}\label{app:fsu}}
\color{black}

\newcounter{filecounter}
\setcounter{filecounter}{1}
\loop
    \edef\filename{\apppath/Приложение. ФСУ/_latex/img/\arabic{filecounter}.pdf}
    \IfFileExists{\filename}{
        \begin{figure}[H]
        \centering
        \includegraphics[width=1\textwidth,height=1\textheight,keepaspectratio]{\filename}
        \end{figure}
        \clearpage
        \stepcounter{filecounter}
    }{\setcounter{filecounter}{0}} % условие выхода из цикла
    \ifnum\value{filecounter}>0
\repeat

\clearpage %!!!!!!!!!!БЕЗ ЭТОЙ СТРОЧКИ ХРЕНЬ



